\appendix{C}{Error-Analysis Examples}

\section{DataBench Error Examples}
\label{sec:databench_error_ex}

\begin{table}[H]
\centering
\small
\begin{tabularx}{\textwidth}{
    >{\raggedright\arraybackslash}p{2.3cm}
    >{\raggedright\arraybackslash}X
}
\hline
\textbf{Query} &
What is the length of the longest text of a review (in chars)? \\
\hline

\textbf{Gold table} &
067\_TripAdvisor (retrieval rank: 11) \\
\hline

\textbf{Gold schema} &
\texttt{ratings, title, text, author, date\_stayed, offering\_id, ...} \\
\hline

\textbf{Alternative tables} &
061\_Disneyland (rank 1): \texttt{Review\_ID, Rating, Year\_Month, Reviewer\_Location, Review\_Text, Branch} \newline
010\_ECommerce (rank 2): \texttt{Clothing ID, Age, Title, Review Text, Rating, ...} \\
\hline

\textbf{Gold answer} &
12592 \\
\hline

\textbf{Prediction} &
20756, derived from 061\_Disneyland \\
\hline

\textbf{Analysis} &
The query does not contain a table-specific cue that uniquely identifies
067\_TripAdvisor. Both 061\_Disneyland and 010\_ECommerce contain explicit
review-text columns and therefore provide plausible interpretations of the
query in the open-domain setting. \\
\hline
\end{tabularx}

\caption{Representative open-domain ambiguity in the DataBench evaluation.
The original closed-domain query does not uniquely identify its designated
gold table after multiple tables become available.}
\label{tab:databench_ambiguity_example}
\end{table}

\section{Financial Error Examples}
\label{sec:financial_error_ex}

